\chapter{Election, Log Repair, and Replica Readmission}
This chapter connects replica state changes into a controlled failover. Keep the model's limitations in mind: the course's authority is a trusted metadata/HW control plane in a single JVM, and leader changes are triggered explicitly by the course; this is not KRaft controller consensus and does not cover dual-leader arbitration during network partitions. Followers read leader data over real TCP, and each broker still writes its own independent log. The course chooses a clean candidate rather than trading data loss for availability; if there is no eligible candidate, the service must remain unavailable.

\begin{figure}[ht]
\centering
\begin{tikzpicture}[font=\small,node distance=12mm]
\node[draw,rounded corners,minimum width=26mm] (stop) {leader A offline};
\node[draw,rounded corners,minimum width=28mm,right=of stop] (elect) {elect B, epoch+1};
\node[draw,rounded corners,minimum width=31mm,right=of elect] (repair) {reconcile log tail};
\node[draw,rounded corners,minimum width=27mm,right=of repair] (admit) {admit replica};
\node[draw,rounded corners,minimum width=27mm,below=of admit] (write) {ALL produce};
\draw[-{Stealth},thick] (stop)--(elect); \draw[-{Stealth},thick] (elect)--(repair);
\draw[-{Stealth},thick] (repair)--(admit); \draw[-{Stealth},thick] (admit)--(write);
\node[below=3mm of repair,align=center] {old epoch rejected\\HW and below cannot change};
\end{tikzpicture}
\caption{A role change must first fence the old role, then repair and prove the replica complete, and only then readmit it to ISR.}
\end{figure}

\lessonsection{25}{Clean Election and Leader Epochs}
\goals{When the old leader is offline, choose only an online replica with the committed prefix and make requests from the old role obsolete.}{Understand eligibility, reported LEO, the HW boundary, and epoch fencing.}
\background{Prerequisite: \stepref{24}{Concurrent Replication} and \stepref{22}{ISR and HW}. A leader is the broker currently allowed to accept writes; a follower has its own partition log. An election selects which known copy continues to serve after a leader failure—it is not an arbitrary choice of any online node. A leader epoch is a number for a role term. It increases when the role changes; requests carry an epoch so a broker can recognize an old term. Epoch fencing rejects later requests from an old role, but cannot undo earlier disk or business side effects. The course authority is a trusted, single-JVM control plane and election is triggered explicitly. If no safe candidate exists, the model prefers unavailability to an unclean election.}
\providededit{Provided \pathcode{ClusterAuthority} (\pathcode{provided/src/main/java/io/simplekafka/cluster/ClusterAuthority.java}) supplies \method{public boolean isOnline(int brokerId)}, \method{public PartitionMetadata metadata(TopicPartition tp)}, and \method{public ReplicationSnapshot snapshot(TopicPartition tp)}; \pathcode{ReplicationSnapshot} is \pathcode{provided/src/main/java/io/simplekafka/cluster/ReplicationSnapshot.java}. The three broker networks and local logs are provided.}{\pathcode{exercises/src/main/java/io/simplekafka/replication/LeaderElection.java}: \method{public PartitionMetadata elect(TopicPartition tp)} and \method{public void checkLeader(int brokerId,TopicPartition tp,int epoch)}.}
\paragraph{Inputs, results, and state.}\pathcode{elect} reads the current leader, epoch, old ISR, online state, reported LEO, and HW for a known topic-partition; it has no separate transport input. If the leader remains online, it returns unchanged metadata without mutating state. A successful change updates authority roles/progress, but does not directly change any broker's disk. \pathcode{checkLeader} checks broker id, partition, and request epoch.
\lessoncontract{A candidate must belong to the ISR at the start of election, be online now, and have authority-reported LEO\(\ge HW\); a longer-looking physical log cannot replace a report. Choose the smallest broker id among eligible candidates, not the highest LEO. On success, increment epoch exactly once, keep HW, reset ISR to only the new leader, and reset other replicas' authority reported LEO to 0 so they must recover and be admitted before acknowledgments count them. An old-epoch request returns \method{FENCED_EPOCH}; the correct epoch from a broker that is not the online leader returns \method{NOT_LEADER}. Wake waiters still awaiting the former leader's acknowledgment.}
\begin{itemize}
\item If the known leader is online, return current metadata and leave epoch/ISR/HW/progress unchanged. With no candidate, return \method{NO_ELIGIBLE_LEADER}; metadata, epoch, HW, and ISR remain unchanged. An online broker outside the old ISR cannot replace it.
\item A negative broker id or epoch is \method{INVALID_REQUEST}; an unknown partition is \method{UNKNOWN_TOPIC_OR_PARTITION}. Election and validation do not make TCP RPCs.
\end{itemize}
\lessonhints{Apply old ISR, online status, and reported LEO\(\ge HW\) in that order before selecting a candidate.}{Distinguish an old epoch (\method{FENCED_EPOCH}) from the correct epoch at a broker with the wrong role (\method{NOT_LEADER}).}{Try a still-online leader, an ISR replica below HW, an online broker outside ISR, and a state with no eligible candidate; compare authority snapshots before and after.}
\expected{Assume leader 1 is offline, epoch=0, old ISR={1,2,3}, and HW=3. The no-retention log starts at offset 0; the common prefix contains three records with null key, one ASCII value byte each, and fixed timestamps, so each occupies 33 bytes. Broker 2 reports LEO=3; broker 3 reports LEO=5; both are online ISR members. Both meet the HW bound, so the smaller id, broker 2, wins even though broker 3 has a longer log. The result is epoch=1, leader=2, ISR={2}, HW=3. Authority-reported LEO for other replicas resets to 0, but election does not delete their physical logs. Broker 3 then gets \method{NOT_LEADER} with epoch 1; an old epoch 0 request gets \method{FENCED_EPOCH}. If all candidates are offline or below HW, even online broker 4 cannot be added, and election leaves state unchanged.}
\paragraph{Tests and completion.}\method{Step25Test.electsTheLowestEligibleIsrAndFencesOldEpochWithoutLosingCommittedPrefix} checks no-op elections while the leader is online, failover, retained HW, and old-epoch fencing; \method{electionRequiresTheCommittedPrefixBeforeChoosingTheLowestIsr} checks that a smaller id below HW is skipped; \method{electionFencesAnOutstandingAllAckWithoutRollingBackAppend} checks that a pending append is not rolled back. The single-step command isolates Step25; the cumulative command runs Steps1–25. A green single-step run is not chapter completion.
\pitfalls{Choosing the online replica with the largest LEO can select a non-ISR or divergent copy; silently choosing a lagging replica when no candidate qualifies can lose the known committed prefix. Real Kafka's controller, epoch history, and consensus are outside this course model; the course does not promise production HA if the authority is lost or the network is partitioned.}
\lessoncommand{25}
\lessonquestions{What minimum safety condition does candidate reportedLEO\(\ge HW\) provide?}{Why can a broker with the correct epoch but not the leader still reject produce?}
\paragraph{Self-check answers.}The condition says that the candidate has at least the prefix the authority marked committed; it does not prove that an uncommitted tail is the same. Epoch identifies the term, while leader identity determines current write permission; both must hold.

\lessonsection{26}{Divergent-Log Reconciliation}
\goals{Compare old and new leader contents at matching offsets, protect their committed prefix, and repair only a suffix allowed to change.}{Understand full-prefix scans, recovery reads, the strict HW boundary, and how long a proof remains valid.}
\background{Prerequisites: \stepref{25}{Clean Election} and the log/truncation and follower-fetch ideas from \stepref{8}{Tail Truncation} and \stepref{21}{Follower Pull Replication}. Equal LEOs mean equal record counts, not equal contents; repair compares each \pathcode{RecordData} at the same offset. The course uses a simple but slower full-prefix scan from offset 0 to the captured leader LEO. Recovery reads can inspect the leader's tail above HW but do not report replication progress. HW is an exclusive boundary: a conflict at offset < captured HW is committed and cannot be overwritten; the offset exactly equal to HW is in the uncommitted tail and may be repaired. A possible truncation point must not be below either captured HW or the current HW rechecked at completion.}
\providededit{Provided \pathcode{RpcClient.call(Messages.Request):Messages.Reply} (\pathcode{provided/src/main/java/io/simplekafka/transport/RpcClient.java}), \pathcode{Messages.ReplicaFetchRequest/ReplicaFetchBody} (\pathcode{provided/src/main/java/io/simplekafka/protocol/Messages.java}), \pathcode{RecoveryProof} (\pathcode{provided/src/main/java/io/simplekafka/cluster/RecoveryProof.java}), and \pathcode{ClusterAuthority} epoch/HW snapshots and partition locks are available; repair does not automatically add the replica to ISR.}{\pathcode{exercises/src/main/java/io/simplekafka/replication/ReplicaReconciler.java}: \method{public long reconcile(int expectedEpoch)}; supplied \method{public Optional<RecoveryProof> recoveryProof()}. Recovery uses \pathcode{REPLICA_FETCH(recoveryRead=true)}, at most 4 records and 4,200,000 bytes per page.}
\paragraph{Inputs, outputs, and state.}\pathcode{expectedEpoch} must match the current epoch. Broker id, topic-partition, local log, borrowed RpcClient, and authority are fixed at construction. Reconciliation captures the leader's reported LEO/HW and the local mutationVersion at offset 0, then makes RPCs outside the lock. On success it returns the local end offset and publishes a proof; recovery reads themselves do not change ISR, reportedLEO, or HW. Entry clears any old proof; a failed validation publishes no new proof.
\lessoncontract{Scan the leader's captured prefix and the local complete prefix to find the first actual content difference or the point where one ends.}
\begin{itemize}
\item A stale expected epoch is \method{FENCED_EPOCH}; an unknown partition is \method{UNKNOWN_TOPIC_OR_PARTITION}; the local broker being the leader or unassigned is \method{INVALID_REQUEST}. A nonzero local retained start, unreadable/noncontiguous leader prefix, committed conflict, or truncation below the final HW is \method{CORRUPT_RECORD}.
\item A differing local record below captured HW is \method{CORRUPT_RECORD} and the file remains unchanged. A missing local record, including a missing committed prefix, can be copied from the leader. A first mismatch exactly at captured HW can be rewritten. A local uncommitted suffix beyond the leader LEO can be truncated.
\item RPC timeout or a local log API mutation during the scan (even a truncate-plus-append returning LEO to the same value) is \method{REQUEST_TIMEOUT}; the log is not repaired and no proof is published. Epoch or leader changes during the scan are \method{FENCED_EPOCH}. Remote error codes propagate; an unexpected body is \method{INVALID_REQUEST}; local disk failure remains \method{STORAGE_ERROR}. Multi-file truncate/append is not guaranteed to roll back atomically on I/O failure.
\item Under authority-lock then local-log-monitor order, recheck epoch, leader, HW, start, LEO, and mutationVersion before repair and proof publication; hold the log monitor through both. The proof binds captured epoch, leader LEO and HW, local log object identity, and final mutationVersion. It is not persistent and does not detect external disk changes that bypass the API.
\end{itemize}
\lessonhints{Separate “the local record is missing” from “a local record exists with different content”; they are different repair situations.}{Treat HW as an exclusive boundary: offsets below it are protected, while the offset equal to it can be the first repairable one.}{Construct a common committed prefix plus a divergent uncommitted tail, then a short/empty replica; check both disk and proof state.}
\expected{Assume offsets 0–2 are common committed records, key=null and one ASCII value byte per record, 33 bytes each; HW=3 and captured leader LEO=4. The old replica has X at offset 3 and the new leader has Y at offset 3, with the same timestamp. Since the first difference is exactly at 3=HW, repair leaves offsets 0–2 unchanged, writes Y at 3, returns LEO=4, and creates a proof; authority ISR/HW do not change. In an independent state where offset 2 differs (2<HW), reconciliation returns \method{CORRUPT_RECORD}, leaves local file bytes and authority unchanged, and has cleared any old proof. An empty local log can instead be filled with offsets 0–3.}
\paragraph{Tests and completion.}\method{Step26Test.repairsOnlyDivergentUncommittedTailAtHighWatermarkAndReturnsLeaderPrefixProof} checks repair at offset=HW and proof publication; \method{refusesMismatchAtHighWatermarkMinusOneWithoutChangingTheLogOrAuthorityAndClearsOldProof} checks rejection at HW-1; \method{fillsShortAndEmptyReplicasWithTheCommittedPrefix}, \method{trimsOnlyAnUncommittedTailWhenTheLocalReplicaHasTheCompleteLeaderPrefix}, and \method{repairsTenRecordsInFourRecordFetchRoundsWithoutReportingRecoveryProgress} check missing/excess tails and four-record pages. Other tests exercise RPC timeout, malformed replies, same-LEO mutations, and epoch/HW races. The single-step command isolates Step26; the cumulative command runs Steps1–26.
\pitfalls{A local HW=0 does not make an existing valid local prefix safe to delete; a missing record is not a content conflict. Real Kafka uses leader-epoch history to locate divergence more efficiently; this course's O(n) scan is a teaching model, not a production repair algorithm.}
\lessoncommand{26}
\lessonquestions{Why handle a conflicting committed record differently from a missing committed record?}{If the leader appends after the scan snapshot, why is the old proof no longer safe for admission?}
\paragraph{Self-check answers.}A conflict means the replicas disagree about a protected committed value, so silently overwriting it is unsafe; a missing record means the local replica can catch up from the available leader prefix. A leader LEO change means the proof did not examine the current tail, so admission must require a new scan.

\lessonsection{27}{Metadata Routing Refresh}
\goals{Route a partition request through cached leader metadata without automatically replaying a business operation after an error.}{Learn bootstrap discovery, cached routes, uncertain outcomes, and why refreshing the control plane is not a data-request retry.}
\background{Prerequisites: \stepref{26}{Divergent-Log Reconciliation}, and the producer and acknowledgment behavior in \stepref{13}{Batched Producer} and \stepref{23}{Produce Acknowledgments}. A bootstrap endpoint is an entry point for discovering the cluster. Metadata gives the current leader endpoint and epoch for each topic-partition. The cached leader is the data-request destination; authority/metadata requests are control-plane traffic, while produce/fetch are data-plane traffic. An uncertain outcome means the broker may have finished appending, but the reply was lost or acknowledgment failed, so the client cannot tell whether a write happened. Automatically resending ordinary produce can create a second record.}
\providededit{Provided \pathcode{RpcClientFactory.connect(Endpoint):RpcClient} (\pathcode{provided/src/main/java/io/simplekafka/transport/RpcClientFactory.java}), \pathcode{RpcClient.call(Messages.Request):Messages.Reply} (\pathcode{provided/src/main/java/io/simplekafka/transport/RpcClient.java}), \pathcode{Endpoint} (\pathcode{provided/src/main/java/io/simplekafka/model/Endpoint.java}), and \pathcode{MessageCodec.encodeRequest/decodeRequest} (\pathcode{provided/src/main/java/io/simplekafka/protocol/MessageCodec.java}) handle connections/frames; broker metadata is supplied by course servers. The existing single-broker RpcClient constructor path remains available.}{\pathcode{exercises/src/main/java/io/simplekafka/client/MetadataRouter.java}: constructor \method{MetadataRouter(List<Endpoint> bootstrap,RpcClientFactory factory)} requires a nonempty bootstrap list with no null entries and a non-null factory; \method{public void refresh(String topic)} and \method{public Messages.Reply call(TopicPartition tp,Messages.Request request)}. The router caches and owns RpcClients created by its factory; \method{close()} closes those clients. Router client constructor paths are also in exercises: \method{SimpleProducer(MetadataRouter,Partitioner,int)}, \method{SimpleConsumer(MetadataRouter,String)}, and \method{GroupConsumer(MetadataRouter,String,String)}; do not remove the single-broker RpcClient overloads.}
\paragraph{Inputs, outputs, and state.}\pathcode{refresh} requests metadata for one topic using bootstrap endpoints in construction order. A connection/factory failure or \method{REQUEST_TIMEOUT} can fall through to the next endpoint. If every endpoint fails, propagate the last failure instead of fabricating a route. Invalid topics or a reply for a different topic return \method{INVALID_REQUEST}. A failed refresh preserves the last successful cache. In Step27, \pathcode{call} accepts matching ProduceRequest or FetchRequest; Step30 later adds \pathcode{IDEMPOTENT_PRODUCE}.
\lessoncontract{\pathcode{call} sends one business request to the cached leader endpoint with the cached epoch. An unsupported request or mismatch between request partition and \pathcode{tp} is \method{INVALID_REQUEST}; absent topic/partition metadata is \method{UNKNOWN_TOPIC_OR_PARTITION}.}
\begin{itemize}
\item On a received \method{NOT_LEADER} or \method{FENCED_EPOCH} response, the router may make a best-effort metadata refresh so a later independent call can use the new route; it must return the original Reply. A refresh failure does not replace that Reply.
\item Other business errors, RPC timeouts, I/O errors, and storage errors do not automatically refresh or replay produce or fetch. In particular, a timeout may happen after the broker appended; an unknown outcome does not mean “nothing happened.”
\item Refresh only reads metadata; the router does not retry the current business operation. A later request is an explicit caller decision. Using a closed router throws \pathcode{IllegalStateException}.
\end{itemize}
\lessonhints{Draw only one business request for a \pathcode{call}; refreshing metadata is another control-plane request, not a second produce.}{For bootstrap fallback, distinguish connection failures from error replies returned by a broker.}{Use an old route to elicit \method{FENCED_EPOCH} or \method{NOT_LEADER}, then compare the returned reply, refreshed cache, and next explicit request.}
\expected{Assume bootstrap order [broker1,broker2,broker3], a cache pointing to broker1/epoch0, and an empty partition. Each example record has null key, a one-byte ASCII value, and a fixed nonnegative timestamp (33 bytes). Broker1 is offline and broker2 has been elected at epoch1. The stale-route request receives an original \method{NOT_LEADER} or \method{FENCED_EPOCH} reply; that same reply returns to the caller while refresh changes the cache to broker2/epoch1. No produce is sent again to broker2 by that call, so it creates no receipt. A later explicit request can receive [0,1), leaving one new log record. If a request instead times out after append, the router still does not replay it; Step29 observes this uncertain result.}
\paragraph{Tests and completion.}\method{Step27Test.refreshSkipsUnavailableAndFactoryThrowingBootstrapsBeforeCachingMetadata} checks bootstrap order; \method{failedRefreshTimesOutWithoutReplacingCacheAndLaterSuccessfulRefreshReplacesIt} checks cache preservation; \method{returnsLeaderErrorsUnchangedAndRefreshesOnlyForTheNextRequest} checks original reply and next-call routing; \method{routerDoesNotReplayAProduceAfterItsAppendResponseTimesOut}, \method{makesRoutedProducerOutcomeUnknownAfterEveryAppendError}, and \method{routerDoesNotReplayStorageOrTransportFailures} check no replay. Other tests cover borrowed-router lifecycle and independent per-topic routes. The single-step command isolates Step27; the cumulative command runs Steps1–27.
\pitfalls{Treating \method{NOT_LEADER} as proof that no record was written can duplicate a message; inventing a success receipt after cache refresh is also wrong. Ordinary \method{SimpleProducer} does not retry automatically. Real Kafka clients have more elaborate retry semantics; this course router refreshes routes but never replays the current business request.}
\lessoncommand{27}
\lessonquestions{Why are refreshing metadata and replaying a business request separate actions?}{Why can a client not infer that the server had no side effect from an error code alone?}
\paragraph{Self-check answers.}Metadata says where future requests should go; it cannot provide a result for an earlier request. A timeout can occur after append, and FENCED can also arise while waiting for acknowledgment. Without a success receipt for that specific request, the caller must treat the result as unknown instead of concluding that nothing was written.

\lessonsection{28}{Replica Readmission to ISR}
\goals{Add a repaired replica to ISR only with fresh prefix proof, and connect repair, readmission, and ALL availability.}{Understand assignment, online/follower role, epoch, log identity, and proof version.}
\background{Prerequisites: \stepref{26}{Divergent-Log Reconciliation} and \stepref{25}{Clean Election}. Repair makes a log a candidate copy of the current leader prefix; ISR readmission is a separate control-plane decision. Only after admission does the replica count toward ALL/minISR and HW. Equal LEOs mean equal record counts, not equal contents. The reconciliation proof binds a prefix comparison to broker, partition, epoch, leader LEO, local log identity, and mutationVersion; a later leader append, target-log change, truncation, or retention makes the evidence stale.}
\providededit{Provided \pathcode{ClusterAuthority} (\pathcode{provided/src/main/java/io/simplekafka/cluster/ClusterAuthority.java}) manages roles/ISR/progress; \pathcode{RecoveryProof} (\pathcode{provided/src/main/java/io/simplekafka/cluster/RecoveryProof.java}) records \pathcode{brokerId,tp,epoch,leaderLEO,localLEO,highWatermark,localLog,localMutationVersion}. Replica reopen and authority state containers are provided.}{\pathcode{exercises/src/main/java/io/simplekafka/replication/ReplicaAdmission.java}: \method{public boolean tryAdd(int brokerId,TopicPartition tp,int epoch)}.}
\paragraph{Inputs, outputs, and state.}\pathcode{tryAdd} makes no RPC. Its target broker, partition, and expected current epoch are arguments; the local PartitionLog is bound at construction. true means already in ISR or newly admitted; false means current state/proof does not permit admission. A successful transition adds the target to ISR, initializes its reported LEO and caught-up clock from the leader LEO, recomputes HW, signals waiters, and consumes/clears the proof.
\lessoncontract{For new admission, verify the target is assigned, online, and a follower for the current epoch; the proof's broker/partition/epoch are correct and its leaderLEO equals the current authority-reported leader LEO; the local log is the same bound object, has the same mutationVersion, start=0, and localLEO=leaderLEO, and proof HW is not above leaderLEO. Under authority lock then local-log monitor, do not allow a local-log write between validation and publication of ISR.}
\begin{itemize}
\item A negative broker id/epoch or unassigned broker returns \method{INVALID_REQUEST}; an unknown partition returns \method{UNKNOWN_TOPIC_OR_PARTITION}; local-log read failure returns \method{STORAGE_ERROR}.
\item A stale epoch, offline target, target that is the current leader, or missing/stale proof returns false without ISR admission. If proof contents or local identity/version/start/LEO disagree, clear the proof; reconcile again before retry. Any target-log append/delete/truncate/close or a change in leader LEO invalidates old proof, even if truncate-plus-rewrite restores the same LEO.
\item An already-in-ISR target returns true idempotently only when it is online, is a follower, and the supplied epoch is current; this repeat call leaves HW unchanged. A stale epoch, offline target, or leader target follows the preceding false-return rule. On first admission, set reportedLEO to leader LEO and caught-up time to the current clock. Admission may succeed even if minISR is not met; ALL preflight can still fail.
\end{itemize}
\lessonhints{Treat proof as short-lived evidence for a particular log object and version, not as a permanent Boolean.}{List the current leader checks and the local identity/version/start/LEO checks that admission repeats.}{Verify a no-change admission, then append at the leader or truncate the target and try the old proof; reconcile before the final admission.}
\expected{Assume offset 0 contains one ordinary record with null key, one ASCII value byte, and timestamp 30 (33 bytes); leader1 and replica2 have copied and reported to LEO=1, so HW=1. Leader1 stops; broker2 becomes epoch1 leader with ISR={2}. Restarted broker1 still has physical LEO=1 but is outside ISR. Reconciliation creates proof for broker1's log and leader LEO1. If leader2 first appends another same-sized record at offset1 with timestamp31, current leader LEO becomes2; calling \pathcode{tryAdd(1,tp,1)} using old proof returns false, ISR/HW remain unchanged, and broker1 remains at LEO1. Reconcile again to copy offset1, then admission returns true, ISR={1,2}, reportedLEO=2, and HW can advance to2.}
\paragraph{Tests and completion.}\method{Step28Test.proofBecomesStaleWhenLeaderAppendsAfterReconciliation} verifies that leader append invalidates proof; \method{admissionRequiresCurrentProofAndOnlineAssignedFollowerAndIsIdempotent} covers missing proof, old epoch, offline/unassigned states, and repeated admission; \method{recoveryProofCannotBeUsedWithAnotherLogObjectWithTheSameRecordsAndVersion} and \method{admissionRejectsProofAfterTruncateAndAfterDifferentDataRestoresSameLeo} cover identity and same-LEO rewrite. \method{recoveredProofRestoresIsrForAllAcksAndReplicaLogsConverge} and minISR integration tests check HW wake-up, ALL, and log convergence. The single-step command isolates Step28; the cumulative command runs Steps1–28, the Chapter 7 completion boundary.
\pitfalls{Equal LEO does not prove equal prefixes; checking only the proof's original epoch without current leader LEO misses a newer tail. Real Kafka ISR admission and replication are more complex; course authority/proof behavior is a teaching model, not a production controller design.}
\lessoncommand{28}
\lessonquestions{Why does equal LEO alone not prove that two logs are safe replicas?}{Why must successful admission update progress, recompute HW, and wake ALL waiters?}
\paragraph{Self-check answers.}Equal lengths can hold different records; a proof compares contents and binds to a particular log version. Adding a replica to ISR can increase the minimum LEO among enough replicas, so HW must be recomputed and waiters notified to recheck acknowledgment conditions.
